\chapter*{Prefacio}
\addcontentsline{toc}{chapter}{Prefacio}
\markboth{Prefacio}{Prefacio}

Este libro tiene dos lectores y cuenta dos historias.

\textbf{El primer lector} es una persona que trabaja por su cuenta o lleva una pequeña empresa y quiere dejar de perseguir sus clientes, sus facturas y sus mensajes entre diez aplicaciones. Para ella, la Parte~II es un manual de uso: cada pantalla de REM, cada botón y para qué sirve.

\textbf{El segundo lector} es quien quiere aprender a construir una aplicación real con ayuda de la inteligencia artificial. Para él, la Parte~III cuenta cómo se hizo REM: las decisiones, la base de datos, la parte agéntica, el diseño, la seguridad y los errores que aparecieron por el camino. REM es un caso de estudio completo: un servidor, una interfaz web, una aplicación de escritorio para macOS, un canal de comunicación cifrado con clientes y un asistente con IA local, todo sin dependencias externas.

\section*{Cómo leerlo}
Para \textbf{usar} REM, lee los capítulos \ref{cap:que-es} a \ref{cap:recorrido} y ve después a la pantalla que te interese en la Parte~II; cada capítulo se puede leer por separado. Para \textbf{aprender a construir}, lee la Parte~III con el código fuente al lado: cada concepto cita su archivo. Si das un \textbf{curso}, los ejercicios finales sirven de prácticas.

\section*{Las marcas numeradas}
Las capturas de la Parte~II llevan círculos naranjas con un número, como este: \mk{1}. Cada número remite a la lista que sigue a la figura, donde se explica qué hace ese botón o campo.

\section*{Una aclaración sobre la IA}
REM usa inteligencia artificial en dos sentidos distintos, y conviene no confundirlos:
\begin{enumerate}
\item \textbf{La IA dentro de la aplicación}: un asistente opcional que redacta borradores y contesta preguntas con los datos del negocio. Funciona con un modelo local, en tu equipo, y \emph{propone}: nada se envía ni se emite sin que tú lo apruebes.
\item \textbf{La IA con la que se construyó la aplicación}: un asistente de programación que escribió código bajo la dirección de una persona, con pruebas y revisión. La Parte~III explica ese método, porque es la parte más reutilizable del proyecto.
\end{enumerate}

\vspace{1em}\hfill\textit{Ramon Guillamon}\\
\hfill\textit{learntouseai@gmail.com}
